% ===========================================================================
% Category: Generalized Regression Neural Network
% Generated from ML_Lab_Manual_Completed.md - edit the Markdown and re-run
% scripts/build_latex_manual.py instead of editing this file by hand.
% ===========================================================================

\chapter{Experiment 19: Generalized Regression Neural Network
(GRNN)}\label{experiment-19-generalized-regression-neural-network-grnn}

\textbf{Category:} Non-parametric regression

\section{1. Objective}\label{objective}

Implement GRNN-style kernel regression from scratch and study the effect
of the smoothing parameter σ with cross-validation.

\section{2. Dataset Used}\label{dataset-used}

Real motorcycle impact data
(\passthrough{\lstinline!data/motorcycle.csv!}, R
\passthrough{\lstinline!MASS::mcycle!}): 133 observations of time after
impact (ms) versus head acceleration (g). Time is standardized before
computing kernel distances.

\section{3. Required Libraries}\label{required-libraries}

numpy, pandas, matplotlib, scikit-learn (API base classes, KFold,
metrics)

\section{4. Brief Theory}\label{brief-theory}

GRNN predicts the Nadaraya-Watson conditional mean: ŷ(x) = Σ\_i y\_i
exp(−\textbar\textbar x−x\_i\textbar\textbar²/2σ²) / Σ\_i
exp(−\textbar\textbar x−x\_i\textbar\textbar²/2σ²). Layer structure:
input → pattern (one Gaussian per training point) → summation (S = Σy\_i
p\_i, D = Σp\_i) → output S/D. Training is one-pass (store the data); σ
controls the bias-variance trade-off.

\section{5. Procedure}\label{procedure}

\begin{enumerate}
\def\labelenumi{\arabic{enumi}.}
\tightlist
\item
  Load mcycle; standardize time.
\item
  Implement the GRNN class (\passthrough{\lstinline!fit!} stores data;
  \passthrough{\lstinline!predict!} computes kernel weights).
\item
  Visualize σ = 0.05 / 0.25 / 1.0 (under/over-smoothing).
\item
  Select σ by 5-fold CV over 0.02\ldots1.0.
\item
  Refit on a 75\% split and evaluate on the 25\% held-out set (R²,
  RMSE).
\end{enumerate}

\section{6. Reference Implementation}\label{reference-implementation}

\begin{lstlisting}[language=Python]
kernels = np.exp(-dist_sq / (2.0 * self.sigma ** 2))   # pattern layer
D = np.sum(kernels, axis=1)                            # denominator
S = np.dot(kernels, self.y_train_)                     # numerator
return S / np.where(D < 1e-12, 1e-12, D)
\end{lstlisting}

\section{7. Results}\label{results}

{\def\LTcaptype{none} % do not increment counter
\begin{longtable}[]{@{}ll@{}}
\toprule\noalign{}
Item & Value \\
\midrule\noalign{}
\endhead
\bottomrule\noalign{}
\endlastfoot
CV-selected σ & 0.080 (CV RMSE 25.00 g) \\
Train R² & 0.8073 \\
Test R² & 0.7255 \\
Test RMSE & 22.78 g (acceleration range −134\ldots+75 g) \\
\end{longtable}
}

σ = 0.05 shows local wiggles (variance), σ = 0.25-1.0 underfits the
sharp deceleration peak. The 5-fold CV curve has a clear minimum at σ =
0.08.

\section{8. Discussion and Conclusion}\label{discussion-and-conclusion}

GRNN is instant to train and needs only one hyperparameter. On the
smooth 1D sine benchmark it reaches R² ≈ 0.95, but on real motorcycle
data the test R² is 0.73 --- a single global σ cannot be simultaneously
small at the impact discontinuity and large elsewhere. This is the
classic bias-variance limitation of global-bandwidth kernel regression;
adaptive bandwidths or local polynomial regression would improve it.

\section{9. Answers to Viva Questions}\label{answers-to-viva-questions}

\begin{itemize}
\tightlist
\item
  \textbf{Why is GRNN called non-parametric?} It does not learn a fixed
  set of parameters; the training data itself forms the model
  (memory-based kernel regression).
\item
  \textbf{What does σ control?} The smoothing bandwidth: small σ fits
  noise (high variance), large σ over-smooths (high bias).
\item
  \textbf{How is GRNN related to kernel regression?} It is exactly the
  Nadaraya-Watson Gaussian kernel regression estimator of E{[}y
  \textbar{} x{]}, with one kernel per training observation.
\end{itemize}

\begin{center}\rule{0.5\linewidth}{0.5pt}\end{center}

\section{10. Complete Code Listing}\label{complete-code-listing}

\emph{Full runnable program for this experiment, extracted from
\passthrough{\lstinline!notebooks/11\_generalized\_regression\_neural\_network.ipynb!}
(executed; results above are its actual output).}

\textbf{Listing 19.1 --- Core libraries}

\begin{lstlisting}[language=Python]
# ---- Core libraries ----
import numpy as np
import pandas as pd
import matplotlib.pyplot as plt
import seaborn as sns

from sklearn.base import BaseEstimator, RegressorMixin
from sklearn.model_selection import KFold, train_test_split
from sklearn.preprocessing import StandardScaler
from sklearn.metrics import r2_score, root_mean_squared_error

# Styling setup
plt.style.use('seaborn-v0_8-whitegrid' if 'seaborn-v0_8-whitegrid' in plt.style.available else 'default')
plt.rcParams['font.sans-serif'] = 'DejaVu Sans'
plt.rcParams['figure.figsize'] = (10, 6)
plt.rcParams['figure.dpi'] = 110
np.random.seed(42)
\end{lstlisting}

\textbf{Listing 19.2 --- GRNN class (from scratch)}

\begin{lstlisting}[language=Python]
class GRNN(BaseEstimator, RegressorMixin):
    """
    Generalized Regression Neural Network (Specht, 1991) - vectorized implementation.

    Prediction is the Nadaraya-Watson kernel regression estimate:
        y_hat(x) = sum_i y_i * exp(-||x - x_i||^2 / (2*sigma^2))
                 / sum_i     exp(-||x - x_i||^2 / (2*sigma^2))

    Training is one-pass: the model simply stores the training set (no gradient descent).
    The only hyperparameter is the smoothing spread sigma.
    """

    def __init__(self, sigma=1.0):
        self.sigma = sigma

    def fit(self, X, y):
        # "Training" = memorizing the data (pattern + summation layers)
        self.X_train_ = np.asarray(X, dtype=np.float64)
        self.y_train_ = np.asarray(y, dtype=np.float64).ravel()
        return self

    def predict(self, X):
        X_test = np.asarray(X, dtype=np.float64)

        # ---- Pattern layer: squared Euclidean distance between each test and train point ----
        # Uses the expansion ||a-b||^2 = ||a||^2 + ||b||^2 - 2 a.b for efficiency.
        dist_sq = np.sum(X_test**2, axis=1, keepdims=True) + \
                  np.sum(self.X_train_**2, axis=1, keepdims=True).T - \
                  2 * np.dot(X_test, self.X_train_.T)
        dist_sq = np.maximum(dist_sq, 0.0)  # guard against tiny negative values from rounding

        # ---- Pattern layer activation: Gaussian kernel of each training point ----
        kernels = np.exp(-dist_sq / (2.0 * (self.sigma ** 2)))

        # ---- Summation layer: D-sum (denominator) and S-sum (numerator) ----
        D = np.sum(kernels, axis=1)                 # unweighted sum of activations
        S = np.dot(kernels, self.y_train_)          # activation-weighted sum of targets

        # ---- Output layer: normalized weighted average (prevents 0/0 for isolated queries) ----
        D_safe = np.where(D < 1e-12, 1e-12, D)
        return S / D_safe


print("Custom vectorized GRNN class defined!")
\end{lstlisting}

\textbf{Listing 19.3 --- Load the real motorcycle impact data
(MASS::mcycle)}

\begin{lstlisting}[language=Python]
# ---- Load the real motorcycle impact data (MASS::mcycle) ----
moto = pd.read_csv('../data/motorcycle.csv')
X = StandardScaler().fit_transform(moto['times'].values.reshape(-1, 1))  # time after impact (scaled)
y = moto['accel'].values                                                 # head acceleration (g)
print(f"Motorcycle data: {len(moto)} observations | time {moto.times.min()}-{moto.times.max()} ms | accel {moto.accel.min():.0f} to {moto.accel.max():.0f} g")

X_eval = np.linspace(X.min(), X.max(), 300).reshape(-1, 1)

# Under-smoothing, near-optimal, over-smoothing
sigmas = [0.05, 0.25, 1.0]
colors = ['purple', 'crimson', 'teal']

fig, axes = plt.subplots(1, 3, figsize=(18, 5))

for ax, sig, col in zip(axes, sigmas, colors):
    # Fit a fresh GRNN for each sigma (training is instant)
    grnn_demo = GRNN(sigma=sig).fit(X, y)
    y_pred_demo = grnn_demo.predict(X_eval)

    ax.scatter(X, y, color='navy', s=25, alpha=0.6, label='Observations')
    ax.plot(X_eval, y_pred_demo, color=col, lw=2.5, label=rf'GRNN ($\sigma={sig}$)')

    # Training R² shows the bias-variance behavior for each sigma
    r2 = r2_score(y, grnn_demo.predict(X))
    ax.set_title(rf"Smoothing $\sigma = {sig}$ (Train $R^2 = {r2:.3f}$)", fontweight='bold')
    ax.set_xlabel("Time after impact (standardized)")
    ax.set_ylabel("Head acceleration (g)")
    ax.legend(loc='lower right')

plt.suptitle("GRNN Behavior on Real Motorcycle Impact Data (σ sweep)", fontsize=15, fontweight='bold')
plt.tight_layout()
plt.show()
\end{lstlisting}

\textbf{Listing 19.4 --- 5-fold CV grid search for the optimal smoothing
sigma}

\begin{lstlisting}[language=Python]
# ---- 5-fold CV grid search for the optimal smoothing sigma ----
sigma_candidates = np.linspace(0.02, 1.0, 50)
cv_scores = []
kf = KFold(n_splits=5, shuffle=True, random_state=42)   # fixed folds for a fair comparison

for sig in sigma_candidates:
    fold_rmses = []
    for train_idx, val_idx in kf.split(X):
        # Split the data into this fold's train/validation parts
        X_tr, y_tr = X[train_idx], y[train_idx]
        X_va, y_va = X[val_idx], y[val_idx]

        # Fit GRNN on the training fold and score on the validation fold
        model = GRNN(sigma=sig).fit(X_tr, y_tr)
        preds = model.predict(X_va)
        fold_rmses.append(root_mean_squared_error(y_va, preds))

    cv_scores.append(np.mean(fold_rmses))   # average CV error for this sigma

best_idx = np.argmin(cv_scores)
best_sigma = sigma_candidates[best_idx]
print(f"Optimal Smoothing Parameter σ: {best_sigma:.3f} (CV RMSE: {cv_scores[best_idx]:.4f} g)")

# ---- Plot the CV error curve and mark the optimum ----
plt.figure(figsize=(9, 4.5))
plt.plot(sigma_candidates, cv_scores, 'b-', lw=2.2)
plt.axvline(x=best_sigma, color='crimson', linestyle='--', label=f'Optimal σ = {best_sigma:.3f}')
plt.title("5-Fold Cross-Validation Error vs Smoothing Parameter (σ)", fontsize=13, fontweight='bold')
plt.xlabel("Smoothing Parameter (σ)", fontsize=11)
plt.ylabel("Cross-Validation RMSE (g)", fontsize=11)
plt.legend()
plt.tight_layout()
plt.show()
\end{lstlisting}

\textbf{Listing 19.5 --- Final model: refit with the CV-selected sigma
and evaluate on a held-out split}

\begin{lstlisting}[language=Python]
# ---- Final model: refit with the CV-selected sigma and evaluate on a held-out split ----
X_tr, X_te, y_tr, y_te = train_test_split(X, y, test_size=0.25, random_state=42)
final_model = GRNN(sigma=best_sigma).fit(X_tr, y_tr)
y_hat = final_model.predict(X_te)

print(f"Final GRNN with CV-selected sigma = {best_sigma:.3f}")
print(f" - Train R²:  {r2_score(y_tr, final_model.predict(X_tr)):.4f}")
print(f" - Test R²:   {r2_score(y_te, y_hat):.4f}")
print(f" - Test RMSE: {root_mean_squared_error(y_te, y_hat):.4f} g")
\end{lstlisting}
